\subsection{The exact period normalizer from the adelic gauge calculation}
\label{sec:flag-aa-normalization}

We retain the original \(\mathrm{MassDataTie}\), the original complete row
system and the measure data of \(\mathrm{bkU}\), not a newly chosen
normalization. Write
\[
 s=r+1,\quad m=2s,\quad n=k+2r+1,\quad a=2r+1,\quad
 l=k-1,\quad f=[F:\mathbb Q],\quad D_F=|\operatorname{disc}(F)|.
\]
Here \(r\) is the flag length (the source's \(\mathrm{rS}\)), \(k\) is
the positive middle rank, and \(n\) is the ambient rank (the source's
\(\mathrm{R.r}\)). The CM quadratic extension \(E/F\) is unramified
at finite places. The two typed rational primes are completely split
in \(E\), their product is \(g\), and the marked prime \(T_u\) of a
row is split and prime to \(g\). We use the standing rank range and
all the row, frame and compatibility assumptions of the original
normalization theorem.

At a finite place \(v\), write \(q_v=|\mathcal O_F/\mathfrak p_v|\),
\(d_v\) for a generator of the different, and \(\eta_v=1\) or \(-1\)
according as \(E_v\) is split or a field. Use the abbreviations
\[
 b_j(v)=\prod_{i=1}^j(1-\eta_v^i q_v^{-i}),\qquad
 L_j=L_F(j,\eta^j),\qquad
 A_v(T)=\frac{b_n(v)}{b_k(v)}
 \begin{cases}1-\eta_v^k q_v^{-k},&v\mid T,\\1,&v\nmid T.\end{cases}
\]
Thus \(A_v(T)\) is exactly the source's \(\mathrm{aP}\).
Put \(A_T=\prod_v A_v(T)\), \(A'=\prod_v A_v(T_{u'})\), and
\(A_{\rm typ}=\prod_{v\mid g}A_v(T_u)\). The products converge absolutely:
every nonconstant infinite Euler factor has index at least \(k+1\);
the factors at marked primes are finite products. All are positive.

\paragraph{The finite self-dual volumes.}
The prescribed standard additive character has different-conductor
pairing, by its restriction at infinity and triviality on the
diagonal field. Its annihilator of \(\mathcal O_{F_v}\) is
\(d_v^{-1}\mathcal O_{F_v}\). Finite unramifiedness of \(E/F\) gives the
perfect integral trace pairing on \(\mathcal O_{E_v}\), also in the
split case and at residue characteristic two. Hence its self-dual
additive measure gives volume \(|d_v|_v\) to \(\mathcal O_{E_v}\) and
volume \(|d_v|_v^{1/2}\) to \(\mathcal O_{F_v}\). On the Hermitian
matrices of size \(h\), the specified pairing is
\(\frac12\operatorname{Tr}_{E_v/F_v}\operatorname{tr}(AB)\):
the diagonal has the ordinary \(F_v\) pairing and each off-diagonal
pair has the trace pairing of \(E_v\). The integral matrix volume
is therefore \(|d_v|_v^{h^2/2}\); in particular there is no missing
factor at two. On \(h\)-tuples in the actual self-dual lattice of
rank \(n\) the volume is \(|d_v|_v^{nh}\). The prescribed Gram-fibre
measure consequently assigns to the primitive tuples of integral
Gram matrix \(B\) the volume
\[
          |d_v|_v^{nh-h^2/2}\,\beta_v(B).                 \tag{AN1}
\]
Here \(\beta_v(B)\) is the normalized primitive density proved in the
earlier finite-place calculation. Indeed the Gram derivative at
a primitive tuple is onto on integral Hermitian matrices:
complete the tuple to an integral basis, use self-duality to choose
the duals, and solve the diagonal equations by the surjectivity
of trace. Each primitive solution modulo \(\pi^j\) then has
\(q_v^{2nh-h^2}\) lifts modulo \(\pi^{j+1}\).
The ratio of tube volume to target-ball volume is (AN1). This is
the value of the actual coarea family: the submersion charts give
continuous integrals on compactly supported tests on the
independent-tuples locus, so the almost-everywhere disintegration
and continuity identify the values at the prescribed Gram matrix.
No regularity of the reduction of \(B\) is needed.

We also need the integral orbit, rather than just the primitive
density. Two primitive tuples of equal Gram matrix are carried to
one another by an integral unitary element. Over the residue field
this is Witt extension for the isometry of their spans; it applies
to degenerate subspaces as well. In the split case, identify the two
components of a tuple with an injection and its dual map; an
isometry respecting their pairings extends by completing the image
basis and the kernel of the dual map. Reduction of the integral
unitary group is onto: the equation for the next correction
\(I+\pi^j X\) in its defining Gram equation is solved by the same
surjective Hermitian derivative. Once the images of the two tuples
agree modulo \(\pi^j\), their difference divided by \(\pi^j\)
prescribes on their common reduced span a map \(L\) satisfying
\(h(Lx,y)+h(x,Ly)=0\) there. Extend this to a skew-adjoint
endomorphism of the full reduced space. To do so, choose a basis
extending the independent tuple, prescribe the entries of the
skew pairing \(h(Lx,y)\) when the first argument is in that tuple,
fill in the opposite entries by skew symmetry, and set the
remaining block to zero. Nondegeneracy gives the endomorphism.
The split case uses the two dual components of this construction.
Lift its infinitesimal correction to an exact isometry by successive
Gram corrections starting at precision \(\pi^{j+1}\). The resulting
infinite product converges integrally and carries the tuples
exactly. This proves the required transitivity.

For the source's ordinary places put \(\alpha_i=d_v^{i-r}\),
\(0\le i\le r\). The upper-entry condition on \(x_{ij}\), \(i<j\),
is \(\alpha_i^{-1}\alpha_j x_{ij}\in\mathcal O_{E_v}\).
The box for the two upper matrices thus has self-dual volume
\[
       |d_v|_v^{2\sum_{i<j}(1-2(j-i))}.
\]
The \(\mathrm{ZZOrd}\) box has entries
\(Z_{ij}\in d_v^{2r-i-j}\mathcal O_{E_v}\), hence volume
\(|d_v|_v^{s^2(2r+1)}\). At a typed place \(d_v\) is a unit.
The upper pair has matching residues, giving volume
\(q_v^{-2\binom{s}{2}}\). In \(\mathrm{ZZTyp}\) write
\(Z=T_F+g_v W\), where \(g_v\) is the typed rational prime
and \(T_F,T_F'\) are the two prescribed Gram matrices.
The remaining condition is
\(W+W^*\equiv (T_F'-T_F)/g_v\pmod {g_v}\).
The right hand side is integral Hermitian, and the reduction map
\(W\mapsto W+W^*\) is onto (including at two). The resulting volume
is \(q_v^{-3s^2}\). These are the original boxes, not estimates
using larger neighborhoods. Their ratios are
\[
 \frac{U_v}{Z_v}=
 \begin{cases}
 |d_v|_v^{-B_r},& v\nmid g,\\
 q_v^{s(2r+3)},&v\mid g,
 \end{cases}
 \qquad B_r=\frac{(r+1)(8r^2+10r+3)}3 .
\]
Since \(\prod_v|d_v|_v=D_F^{-1}\), it follows that
\[
    \frac{\mathrm{unitLevelVolume}}{\mathrm{volZZ}}
             =D_F^{B_r}\,g^{s(2r+3)f}.                  \tag{AN2}
\]

At a typed place, in the actual integral hyperbolic splitting,
the graded flag stabilizer has unipotent radical of dimension
\(d_N=r(2k+2r+1)\) over \(F_v\), with the fixed-last-vector
middle Levi of rank \(k\). Over the split residue algebra this
Levi is \(\mathrm{GL}_k(\mathbb F_{q_v})\). All its residue
elements lift, as do the unipotent parameters. This can be seen
in the upper, middle and skew coordinates: use arbitrary integral
upper and middle entries and solve the diagonal trace equations
integrally. Its reduction consequently has order
\(q_v^{d_N+k^2}b_k(v)\). The full integral group reduces onto
\(\mathrm{GL}_n(\mathbb F_{q_v})\), of order \(q_v^{n^2}b_n(v)\).
The original \(K_{\rm nat,v}\) is exactly the inverse image of
the first image. Since \(n^2-k^2-d_N=s(2k+2r+1)\), this proves
the exact index, and hence
\[
        \frac{\mu K_{\rm nat}}{\mu K}
          =\frac{g^{-s(2k+2r+1)f}}{A_{\rm typ}}.         \tag{AN3}
\]
For clarity, the local volume of a pair of primitive tuples
used in \(\mathrm{Mp}\) has the same index, not the index
of a pointwise stabilizer alone. If \(S_0\) fixes the first
\(s\)-tuple pointwise and \(S\) fixes its flag and graded
vectors, then \(S/S_0\) is the upper unitriangular group of
size \(s\). Its integral image is onto the corresponding
upper-entry box: the hyperbolic splitting and primitivity of
the last vector permit every integral shift of that vector.
The exact local compatibility in \(\mathrm{MassDataTie}\)
disintegrates \(\mu_S\) as the pointwise Haar measure and
this self-dual upper-entry measure. Therefore the volume of
one full primitive orbit is the single upper-box volume
times \(\mu K/\mu S\). At the typed places the two full
tuples must have identical residues; for each first residue
the proportion of the second orbit in that residue is
the reciprocal of the number of primitive residues of the
fixed Gram matrix. The residue orbits are transitive by Witt
extension, and congruence of the two row norms makes their
residue Gram matrices equal. The pointwise-stabilizer image
has index \(q_v^{2\binom{s}{2}}\) in the graded-stabilizer
image, precisely the inverse proportion for the upper pair.
Combining the two proportions gives
\[
   \mathrm{Mp}_v
       =U_v\,\frac{\mu K_v\mu K_{{\rm nat},v}}
                        {\mu S_v\mu S'_v}\quad(v\mid g).
\]
The ordinary product in the concrete datum is set equal to
the corresponding ordinary expression by definition
(\(\mathrm{MassDataConc.bC}\)). Consequently
\[
       \mathrm{MM}=U\,\frac{\mu K\mu K_{\rm nat}}{\mu S\mu S'},
       \qquad U=\mathrm{unitLevelVolume}.              \tag{AN4}
\]
In particular the calculation is valid with the actual Haar
normalizations, without assuming they are probability measures.

The complete standard tuple \(\mathrm{tupH}\) has \(a=2r+1\)
vectors and Gram matrix the \(r\) hyperbolic pairs and the marked
line of norm \(-T\). Locally the two hyperbolic bases may be
rescaled by the reciprocal different-ideal generators to become
integral, with determinant of the tuple change equal to one.
This changes neither the Gram determinant nor the fibre volume.
The resulting tuple is primitive in the original split frame
lattice. Its residue Gram corank is zero off \(T\) and one
at \(T\). By (AN1), integral orbit transitivity, and the exact
primitive-density formula, the compact orbit of the complete
\(u'\)-tuple has volume
\[
    \Omega_v=|d_v|_v^{na-a^2/2} A_v(T_{u'}),\qquad
    \Omega:=\prod_v\Omega_v=D_F^{-(na-a^2/2)}A'.         \tag{AN5}
\]
The fixed-last-vector row mass is the exact mass of
Lemma~\ref{lem:flag-row-genus-mass}; writing \(C(E,k)\) for the
unimodular Hermitian mass constant used in (GM1), it is
\[
 \rho_0=R_T=C(E,k)T^{kf}P_T,\qquad
 P_T=\prod_{v\mid T}(1-\eta_v^kq_v^{-k}),\qquad T=T_u.    \tag{AN6}
\]
Indeed the cyclic discriminant module of the rank-\(k\) complement
has \((T-1)^f\) gluing choices at the split marked places. They
cancel the denominator of its primitive-last-vector mass. Finally
the definitions of the normalized local factors give
\[
 \mathrm{al}=\prod_v\mathrm{alpD}_v=\frac{A_T A'}{A_{\rm typ}},
 \qquad
 A_T\prod_{i=k+1}^n L_i=P_T.                            \tag{AN7}
\]

\paragraph{The global quotient-volume identity.}
This step can use a classical theorem instead of the full
regularized Siegel--Weil formula. For a nondegenerate independent
rational tuple \(X_0\) with Hermitian Gram matrix \(B_0\), let \(G\)
be the ambient unitary group and \(H\) its pointwise stabilizer
(the unitary group of the orthogonal complement). Both have
positive rank here, both have determinant character module with
the same quadratic Galois action, and both have Tamagawa number
\(2\). We use the standard unitary Tamagawa measure normalization,
as explicitly recorded in Harris--Li--Skinner,
\emph{p-Adic L-functions for unitary Shimura varieties, I},
Documenta Math. Extra Volume Coates (2006), 393--464, (0.2.1).
The equality \(2\) for all Hermitian forms also follows from the
classical Tamagawa theorem for simply connected special unitary
groups, Ono's formula for the quadratic norm-one torus, and the
Hasse norm theorem; the determinant map has an algebraic section
along any nonisotropic rational line. Equivalently it is the
unitary case of Kottwitz's Tamagawa-number formula.

The Gram map is a submersion on the independent-tuples locus.
A rational invariant fibre gauge \(\omega_{\rm fib}\) is the ratio
of rational constant additive top forms in tuple and Hermitian
matrix coordinates. It is invariant since the unitary action
preserves the absolute ambient determinant, in fact its
\(F\)-linear determinant is one. Witt extension identifies the
nondegenerate fibre with \(G/H\), locally and rationally on its
points. Choose rational invariant forms on \(G,H\) whose quotient
on that fibre is \(\omega_{\rm fib}\), using the tangent exact
sequence at \(X_0\). The specified self-dual additive gauges differ
from the rational constant coordinate gauges by the square roots
of the absolute determinants of their rational pairings and the
scalar self-dual different factors. Multiplying the determinant
factors over all places gives one by the product formula;
the different factors are exactly the dimensional discriminant
normalization in the Tamagawa measures. The two groups have the
same \(L_F(1,\eta)\) convergence factors (their character modules
agree), so these cancel in their quotient, globally and at each
finite place. Thus the Tamagawa quotient on the orbit is exactly
the product of the given self-dual Gram-fibre measures. At good
finite places the quotient compact orbit volumes are
\(b_n(v)/b_{\dim_E H}(v)\); the product converges absolutely in
the present rank range, so this comparison uses the actual
restricted product, not a rearrangement of a conditionally
convergent Euler product.

By uniqueness of Haar measures, the actual \(\mu_G,\mu_H\)
in \(\mathrm{MeasCompat}\) are proportional to these Tamagawa
measures. The proportionality constants are equal, since their
quotient is prescribed to be the same Gram-fibre measure.
The equality of Tamagawa numbers now proves, for the two
complete standard row tuples and the original nondegenerate
\(m\)-tuple \(X_0\),
\[
      \mathrm{volHa}=\mathrm{volHa}'=\mathrm{volB}=\mathrm{volG}.
                                                               \tag{AN8}
\]
This uses only the stipulated measure compatibility and the
classical Tamagawa theorem. It does not require a Siegel--Weil
identity on the degenerate Gram locus.

For the \(u'\) complete standard tuple, write \(a_\infty>0\)
for the scalar in the infinite Weil quotient, with the
prescribed infinite Gram-fibre measure:
\(d\mu_{G,\infty}=a_\infty\,d\lambda_\infty\,d\mu_{H',\infty}\).
The analogous finite scalar is \(a_\infty^{-1}\) by the
actual product compatibility. One can see this by testing the
global Weil formula on products of nonzero compactly supported
functions at infinity and characteristic functions of compact
opens at the finite adeles; uniqueness of the local invariant
quotient measures gives the two factors and forces their
product to be one. The compact orbit at finite adeles is
precisely \(K/(K\cap H')\), with measure \(\Omega\) by (AN5).
It follows that
\[
                    \mu K\,a_\infty=\mu H'\,\Omega.     \tag{AN9}
\]
This statement is unchanged if the chosen real and finite
Haar factors are rescaled in opposite directions.

\paragraph{The infinite quotient and the two-sided Jacobian.}
We give the real calculation to specify all constants in
(AN9) and in the disintegration of the original real form.
At a real place choose the actual orthogonal Witt basis:
the \(r\) hyperbolic pairs, a positive orthonormal basis of
rank \(k\), and the normalized negative last vector.
The matrix of the form in these coordinates is Hermitian
unitary. Changes between such coordinates preserve the
Hilbert--Schmidt Lie-gauge measure: identify the Lie algebra
with skew-Hermitian matrices by multiplication by the form
matrix, and note that the induced congruence determinant
is \(|\det P|^{2n}=1\). Thus this computation applies to
the actual coordinate transport, not just to a special word.
Denote by \(c_G\) the scalar of the given real Haar measure
relative to this Lie gauge, by \(h,h'\) the total masses of
the two compact pointwise complete-frame groups, and by
\(\kappa,\kappa'\) the scalars of the two real unipotent
Haar measures in the reference Witt coordinates.

Let \(V_j\) be the Hilbert--Schmidt Lie-gauge volume of
\(\mathrm U(j)\), with \(V_0=1\). The projection to its last
unit vector has sphere fibre \(\mathrm U(j-1)\). In the
horizontal tangent space, each complex off-diagonal entry
contributes twice its squared modulus to the Lie norm;
the imaginary diagonal entry contributes its ordinary squared
modulus. Accordingly its quotient density is \(2^{j-1}\)
times the sphere's Euclidean density. Integration yields
\[
          \frac{V_j}{V_{j-1}}
             =2^{j-1}\frac{2\pi^j}{(j-1)!}
             =\frac{(2\pi)^j}{(j-1)!}.                 \tag{AN10}
\]
This is the source's \(\mathrm{UnitaryVolumeMass.volU}\), not
the volume in the gauge where independent off-diagonal complex
entries have raw Lebesgue measure.

For a fixed nondegenerate real \(a\)-tuple of complete-frame
Gram determinant of modulus \(T'\), the self-dual Gram-fibre
measure is \(T'^{\,k}\) times the Hilbert--Schmidt quotient
gauge on \(G/H'\). To check this determinant explicitly,
write the frame and orthogonal complement as the columns
of a matrix \(P\), so \(P^*JP=B\perp I_k\) and
\(|\det P|^2=|\det B|=T'\). In the tuple tangent coordinates
\(dX=P\binom{U}{Z}\) use \(S=BU\). The ambient real Lebesgue
Jacobian is \(|\det B|^a\), and the inverse substitution
\(U=B^{-1}S\) contributes \(|\det B|^{-2a}\).
The Gram derivative is \(S+S^*\): the diagonal real part
has derivative 2 and each complex upper target can be
used directly, leaving independent skew coordinates.
The self-dual tuple gauge is \(2^{na}\) times coordinate
Lebesgue measure and the self-dual Hermitian target gauge is
\(2^{a(a-1)/2}\) times its coordinate measure. The quotient
therefore has density
\(2^{na-a(a+1)/2}|\det B|^{-a}\,dS_{\rm skew}\,dZ\).
On the group, skew-matrix congruence by \(P\) has Jacobian
\(|\det P|^{2n}=|\det B|^n\); removing the lower-right
\(\mathfrak u(k)\) coordinates leaves density
\(2^{na-a(a+1)/2}|\det B|^{-n}\,dS_{\rm skew}\,dZ\).
Their ratio is \(|\det B|^{n-a}=T'^k\), as claimed.
Both gauges are invariant, so the tangent computation
determines the full quotient measure. Thus the real scalar
in (AN9) is
\[
              a_v=\frac{c_G V_k}{h'T'^k},\qquad
              a_\infty=\prod_{v\mid\infty}a_v.          \tag{AN11}
\]

For the two-sided real disintegration, first separate the
full flag unipotent groups on the big cell. The conditions
that all the leading cross-pairing minors through rank \(r\)
are nonzero are exactly the open elimination conditions.
Triangular row and column operations make their cross
pairing diagonal, with entries \(c_1,\ldots,c_r\).
The radical parameters on each side then put the images
of the isotropic vectors and their duals in the opposite
standard positions: their middle projections determine
the middle parameters, and isotropy fixes the symmetric
part of the remaining top projections. The independent
skew parameters complete that elimination. The remaining
isometry is on the \(k+1\) dimensional middle space.
This construction is unique: after the triangular
elimination, an ambiguity between the two radical
elements would fix both opposite full isotropic frames
and act trivially on the middle quotient, hence is the
identity. The middle element is then determined.
In these coordinates the big cell is
\(n_1 w(c) b n_2^{-1}\) with \(b\in\mathrm U(k,1)\).
The opposite torus \(w(c)\) exchanges the hyperbolic
pairs with coefficients \(c_i\) and \(\overline c_i^{-1}\),
up to the harmless signs of the prescribed Weyl element.
Its diagonal complex parameter has Lie-gauge density
\(2\,d^2c_i/|c_i|^2\): both dual entries contribute to
the squared norm. Left Maurer--Cartan differentiation
at this section gives the determinant of
\(\operatorname{Ad}(\operatorname{diag}(c,I,\bar c^{-1}))\)
on the positive unipotent. An upper \(ij\) entry has
weight \(c_i/c_j\), the \(k+1\) middle entries for index
\(i\) have weight \(c_i\), and the skew block has weights
\(c_i\bar c_j\) (real weight \(|c_i|^2\) on its diagonal).
The exponent of \(|c_i|\) in the real determinant is
\(2(r-2i+1)+2(k+1)+2r\), using \(1\le i\le r\).
Combining the torus differential with this determinant gives
\[
         2^r\prod_{i=1}^r |c_i|^{2(k-1)+4(r-i+1)}d^2c_i.
                                                               \tag{AN12}
\]
The density is independent of \(n_1,n_2,b\) by unimodularity,
translation invariance, and the middle action of determinant
modulus one on the radical. The Lie subspaces at the Weyl
element are orthogonal, so there is no additional constant.
The complement of the big cell is a proper real algebraic
set in the connected unitary group, hence Haar null.
This is the actual big-cell Haar formula in the coordinates
of \(\mathrm{JacConst}\).

In the terminal group \(\mathrm U(k,1)\), put
\(\zeta=h(z,bz')\) (conjugating this coordinate does not
change its planar measure). In orthogonal coordinates a
vector of norm \(-T'\) has form \((x,y)\) with
\(|y|^2=T'+\|x\|^2\). The one-vector coarea measure for
the prescribed trace pairing is \(2^k d^{2k}x\,d\arg y\),
since the norm constraint has radial derivative \(2|y|\).
As \(|\zeta|^2=T(T'+\|x\|^2)\),
\[
 R^{2k-1}dR
   =T^{-k}(|\zeta|^2-TT')^{k-1}|\zeta|\,d|\zeta|,
        \qquad R=\|x\|.
\]
The compact group fixing \(z\) acts transitively on the
sphere direction, with normalized pushforward of its Haar
measure equal to normalized sphere area. The quotient by
the right compact group has coarea factor \(T'^k\)
by the same tangent computation used in (AN11), now with
one vector. Hence the terminal Haar density, with compact
measures of masses \(h,h'\), is
\[
 \frac{2 V_k^2}{(TT')^k V_{k-1}hh'}
       (|\zeta|^2-TT')^{k-1}d^2\zeta.                  \tag{AN13}
\]
The endpoint \(R=0\) has measure zero. This holds for any
measurable choice of terminal section, by the compact
averages; the section used in the original datum is one
such choice.

In the original real variable the cross Gram matrix is
diagonal with these \(r+1\) entries. Its weight is exactly
\[
     (|\zeta|^2-TT')^{k-1}
               \prod_{i=1}^r |c_i|^{2(k-1)+4(r-i+1)} .
\]
The original \(\mathrm{dcInf}\) is twice planar Lebesgue
measure on each complex coordinate. Therefore the factors
\(2^r\) in (AN12) and 2 in (AN13) are absorbed, not lost.
After including the original Haar scalars, the infinite
disintegration constant at this place is
\[
         C_v=\frac{c_G V_k^2}
               {(TT')^k V_{k-1}hh'\kappa\kappa'} .
\]
Use (AN11), (AN10), and the exact unipotent covolume theorem
\ref{global-gram-haar-covolume}. In its reference
Hilbert--Schmidt gauges the two original infinite unipotent
covolumes are both \(D_F^{e_r/2}\), with
\[
              e_r=-\frac{8r^3+6r^2+6kr^2+r}{3}.
\]
The actual measures multiply these by the respective products
of \(\kappa,\kappa'\). With \(C=\prod C_v\) the original
unscaled \(\mathrm{DisintegrationConstant}\), this gives
\[
 C\,\mathrm{volH}\,\mathrm{volN}\,\mathrm{volN}'
    =a_\infty D_F^{e_r}
          \left(\frac{(2\pi)^k}{(k-1)!}\right)^f T^{-kf}.
                                                               \tag{AN14}
\]
The calculation includes \(r=0\) with empty upper-coordinate
products. It proves the reference Haar bracket in the
normalization theorem without a bound on \(c_G\).

\begin{theorem}[The original unscaled \(AA\) identity]
\label{thm:flag-aa-normalization}
For the actual \(\mathrm{bkU}\) datum under the original
\(\mathrm{MassDataTie}\) and the admitted row and finite
level hypotheses, the equality
\[
 \boxed{\displaystyle
  \mathrm{cExpl}\,\mathrm{massRows}
           \prod_v\mathrm{alpD}_v
              =AA\,g^{mf(k-1)} }                       \tag{AN15}
\]
holds, where \(AA=(\mathrm{bkU}).AA\). All the measures on
both sides are the original measures.
\end{theorem}
\begin{proof}
At finite level the actual graded stabilizer has integral
factorization \(S\cap K=(N\cap K)(H\cap K)\).
Indeed the projection to the self-dual middle quotient of
the standard integral splitting preserves its lattice and
last vector; every such middle isometry has the block
diagonal integral lift, leaving an integral unipotent.
The same argument works for the second word. Product Haar
compatibility consequently gives \(\mu S=\mu N\mu H\) and
\(\mu S'=\mu N'\mu H'\). Unipotent strong approximation
(from the rational triangular coordinates and their successive
additive quotients) identifies the real period fundamental
domain times \(K_N\) with an adelic fundamental domain, so the
stipulated total unipotent covolume 1 gives
\(\mathrm{volN}=1/\mu N\) and \(\mathrm{volN}'=1/\mu N'\).
The compact middle adelic quotient decomposes according to
the original row genus; each class has reciprocal stabilizer
weight, hence
\(\mathrm{volHa}=R_T\,\mathrm{volH}\,\mu H\).
This is the same row mass as (AN6).

The definition of \(AA\) in the concrete datum is
\(\mathrm{volB}\,C\,\mathrm{MM}/(\mathrm{volZZ}\mu K)\):
the factor \(\mathrm{kapN}\) in the scaled \(\mathrm{Coo}\)
has been removed in \(\mathrm{bkU}\).
Substitute (AN4) and the preceding covolume identities:
\[
 \frac{AA}{R_T}
   =\frac{\mathrm{volB}}{\mathrm{volHa}}
     C\,\mathrm{volH}\,\mathrm{volN}\,\mathrm{volN}'
       \frac{U}{Z}\frac{\mu K_{\rm nat}}{\mu H'}.
\]
All denominators are positive by the Haar constructions.
Use (AN8), (AN9), (AN2), (AN3), (AN5), (AN14):
\[
 \frac{AA}{R_T}
   =D_F^{e_r+B_r-(na-a^2/2)}
       \left(\frac{(2\pi)^k}{(k-1)!}\right)^f T^{-kf}
         \frac{A'}{A_{\rm typ}}\,g^{-mf(k-1)} .
\]
The exponent simplifies identically to
\(-2r(r+1)(k-1)-k+\frac12\).
By the exact Hermitian mass formula in (GM1),
\[
  \frac{C(E,k)}{C(E,k-1)}
    =D_F^{k-1/2}\left(\frac{(k-1)!}{(2\pi)^k}\right)^f L_k .
\]
Multiplying by (AN6) and using (AN7) now yields
\[
 AA\,g^{mf(k-1)}
  =C(E,k-1)D_F^{-2r(r+1)(k-1)}
                 \left(\prod_{i=k}^n L_i\right)
                 \frac{A_T A'}{A_{\rm typ}} .
\]
The first three factors are exactly the defining
\(\mathrm{cExpl}\,\mathrm{massRows}\) (including the
different exponent), and the last ratio is exactly
\(\prod_v\mathrm{alpD}_v\). This proves (AN15).
\end{proof}

In particular (GM11) is now supplied, and the generic mass
comparison (GM12) holds under these actual hypotheses:
\[
 \mathrm{ngen}_{\rm tot}\le
 2\,AA\,\rho^{2r(r+1)f(k-1)}H_{\rm tot}^{s(s-1)}\,\mathrm{nPt}.
                                                               \tag{AN16}
\]
The remaining generic local supplier must still relate the
prescribed pinned real quotient measures to the measures
in the class factorization and retain the correct scale
in the positive point sum. This normalization proof alone
does not bound the nonzero-Gram orbital energies.
